% =====================================================================
% chapters/00_Abstract.tex — rendu LaTeX de ../traductions/00_abstract.en_brut.txt (23 septembre 2026)
% Chapitre 0 de la VERSION FRANÇAISE de l'article court : le résumé. \input depuis main.tex,
%   dans le préambule, avant \begin{document}, comme le fait le gabarit AAMAS.
% Rendu mot pour mot de la traduction fournie par l'auteur.
% TYPOGRAPHIE : babel french pose seul l'espace fine devant ; : ? !
%   Ce qui est explicite ici : 1\,000 et 3\,299 (séparateur de milliers en espace fine
%   insécable), la virgule décimale étant un simple caractère.
% NOTE DE BAS DE PAGE : préservée sur l'enquête CEREMA 2023.
% =====================================================================

\begin{abstract}
Les simulations multi-agents de mobilité urbaine choisissent les modes de transport à l'aide de modèles tabulaires ajustés à partir d'enquêtes sur les déplacements déclarés, ou à l'aide de règles élaborées par des experts du domaine. Ces deux approches fixent l'espace des comportements représentables avant le lancement de la simulation ; ainsi, un facteur non pris en compte par une variable ou une règle n'influence aucune décision. Les agents génératifs construits sur des modèles de langage promettent de lever cette contrainte, en intégrant des heuristiques de décision que les enquêtes n'enregistrent jamais. À notre connaissance, la validité de cette promesse face à une population réelle n'a pas encore été vérifiée.

Nous l'évaluons sur la région toulousaine, en nous appuyant sur l'enquête certifiée CEREMA 2023 sur les déplacements des ménages\footnote{Enquête Ménages Déplacements (EMD), Toulouse / Grande agglomération toulousaine --- 2023, CEREMA, Syndicat mixte des transports en commun de l'agglomération toulousaine (producteurs), PROGEDO-ADISP (diffuseur). Rapport final : \url{https://www.aua-toulouse.org/wp-content/uploads/2024/05/Rapport-final-68-pages-Enquete-mobilite-2023-Bassin-de-vie-toulousain.pdf}}. Une cohorte fermée de 1\,000 profils synthétiques, contrôlés par rapport à treize marges de cette enquête, effectue les 3\,299 déplacements du jour que nous analysons. Quinze décideurs effectuent ces déplacements selon un protocole unique. Ce protocole leur fournit les mêmes 21 variables d'entrée. Il calcule la probabilité que chaque décideur attribue à chaque option proposée.

En ajustant l'invite pour mettre en évidence des critères généraux (tels que le confort, les contraintes ou les opportunités liées au choix d'une option), l'agent génératif se rapproche des caractéristiques des modèles tabulaires. Aucun de nos agents ne les atteint. Un classificateur qui lit le même contexte avec un type de sortie fixe et n'écrit aucun texte, les égale pour un cinquantième du coût. L'agrégat masque ce qu'une lecture au niveau de la décision révèle. Deux décideurs que l'agrégat ne peut distinguer sont en désaccord sur un trajet sur trois. Le classificateur se trompe sur plus de trajets individuels qu'une règle qui ne lit rien.

La délibération verbalisée trouve sa place en dehors du quotidien. Chez un agent, une panne de moteur est entrée dans la mémoire, a dissuadé l'utilisation de la voiture pendant quinze jours, puis s'est estompée. Aucun décideur ne disposant que des 21 variables ne peut répondre à un tel événement. Nos mesures encouragent l'utilisation du modèle de langage lorsque les événements surviennent et laissent le quotidien à un décideur qui ne délibère pas.
\end{abstract}
